\scanpage{6}
% Reused typescript, read after rotating the supplied image by 180 degrees.
\mspage{5}
\begin{center}-3.1-\end{center}
\uline{Chapter} \uline{3}. \uline{An} \uline{elementary} \uline{example}

The purpose of this chapter is to give, independently of the
forgoing theory, an exhaustive analysis of a single special class of
actions. This has two objectives: first of all, the concepts of
pre-stable and stable are worked out in a representative non-trivial
case. Secondly, by attacking this case directly, we can circumvent
the difficulties involved in extending the previous work to semi-simple
groups in characteristic $p$ (cf. Preface). And, in fact, we obtain the
result over the ring of integers, hence, \uline{a} \uline{fortiori}, over any field.
This will enable us later to construct the so-called arithmetic
schemes of moduli of abelian varieties.

We have no use for a ground field in this chapter: all the
schemes that we will consider will be of finite type over $\operatorname{Spec}(\mathbb Z)$.
$\mathbb A^n$, $\mathbb P_n$, and $\operatorname{PGL}(n)$ stand for the usual schemes \uline{over} $\uline{\mathbb Z}$. The
action we shall study is the canonical action:
\[
\sigma_m:\operatorname{PGL}(n)\times(\mathbb P_n)^{m+1}\longrightarrow(\mathbb P_n)^{m+1}
\]
\margin{$-$}
i.e. of projective transformations on sequences of points in projective
space.

\fbox{§1.} The first step is to obtain a large number of\inserted{\ }invariant affine
open subsets in $(\mathbb P_n)^{m+1}$. To this end, let $p_0,\ldots,p_m$ denote
the projections of $(\mathbb P_n)^{m+1}$ onto its $m+1$ factors. Let
$L_i=p_i^*(\underline{o}_{\mathbb P}(1))$. Now unfortunately, $\underline{o}_{\mathbb P}(1)$ admits no $\operatorname{PGL}(n)$-
linearization with respect to the action of $\operatorname{PGL}(n)$ on $\mathbb P_n$. However, if
